\documentclass{article}
\usepackage[utf8]{inputenc}
\usepackage{a4wide}

\begin{document}

\section*{เมืองในฝัน}

อาณาจักรในฝันของจักรินมีโครงสร้างเป็นต้นไม้ (Tree) โดยมีเมืองทั้งหมด n เมือง

แต่ละเมืองมีการตั้งลำดับตั้งแต่ 1 ถึง n โดยเมืองหลวงจะเป็นลำดับที่ 1

ในอาณาจักรมีถนนที่เดินทางได้สองทิศทางระหว่างเมืองจำนวน n-1 เส้น

ซึ่งทำให้สามารถเดินทางจากเมืองหนึ่งไปยังเมืองอื่นได้โดยการเดินทางผ่านถนนต่าง ๆ ในต้นไม้



จักรินอยากให้อาณาจักรในฝันเป็นเมืองที่มีความเจริญก้าวหน้า

และยังคงรักษาสมดุลของแหล่งท่องเที่ยวต่างๆ ตามธรรมชาติ

จักรินจึงมีนโยบายว่าในอาณาจักรนี้เมืองทุกเมืองสามารถถูกเลือกให้พัฒนาเป็นเมืองเน้นอุตสาหกรรมหรือเมืองเน้นท่องเที่ยวเชิงธรรมชาติ

จักรินต้องการเลือก k เมืองที่จะพัฒนาเป็นเมืองเน้นอุตสาหกรรม โดยเมืองอื่น ๆ

จะพัฒนาเป็นเมืองเน้นท่องเที่ยวเชิงธรรมชาติ

เมืองหลวงเองก็ต้องถูกพัฒนาเป็นเมืองเน้นอุตสาหกรรม

หรือเมืองเน้นท่องเที่ยวเชิงธรรมชาติอย่างใดอย่างหนึ่งเช่นกัน ทุกปีจะมีการประชุมที่เมืองหลวง

โดยทุกเมืองที่เน้นอุตสาหกรรมจะต้องส่งทูตไปที่เมืองหลวง

ทูตจะเดินทางไปยังเมืองหลวงด้วยเส้นทางที่ใกล้ที่สุด,

โดยค่าความสุขของทูตจะเท่ากับจำนวนเมืองเน้นท่องเที่ยวเชิงธรรมชาติ

ที่เขาเดินทางผ่านบนเส้นทางมาเมืองหลวง



\ul{งานของคุณ}: ช่วยจักรินเลือกเมืองเน้นอุตสาหกรรมจำนวน k เมือง

เพื่อให้ผลรวมของความสุขของทูตทั้งหมดมากที่สุด



\hfill\break



{\def\LTcaptype{none} % do not increment counter

\begin{longtable}[]{@{}

  >{\raggedright\arraybackslash}p{(\linewidth - 4\tabcolsep) * \real{0.3333}}

  >{\raggedright\arraybackslash}p{(\linewidth - 4\tabcolsep) * \real{0.3333}}

  >{\raggedright\arraybackslash}p{(\linewidth - 4\tabcolsep) * \real{0.3333}}@{}}

\toprule\noalign{}

\begin{minipage}[b]{\linewidth}\centering

ข้อมูลเข้า

\end{minipage} & \begin{minipage}[b]{\linewidth}\centering

ข้อมูลออก

\end{minipage} & \begin{minipage}[b]{\linewidth}\centering

คำอธิบาย

\end{minipage} \\

\midrule\noalign{}

\endhead

\bottomrule\noalign{}

\endlastfoot

\begin{minipage}[t]{\linewidth}\raggedright

7 4\\

1 2\\

1 3\\

1 4\\

3 5\\

3 6\\

4 7\\

\strut

\end{minipage} & 7 & \begin{minipage}[t]{\linewidth}\raggedright

ในกรณีนี้ เมือง 1 เป็นเมืองหลวง\\

(สามารถเลือกเป็นทั้งเมืองอุตสาหกรรมหรือเมืองท่องเที่ยว)\\

และจักรินต้องเลือก 4 เมืองที่จะเป็นเมืองเน้นอุตสาหกรรม จะได้ตามรูป\\

\strut \\

\includegraphics[width=2.11458in,height=1.71875in,alt={image.png}]{/public/upload/d1107ef906.png}\\

\strut

\end{minipage} \\

\begin{minipage}[t]{\linewidth}\raggedright

4 1\\

1 2\\

1 3\\

2 4\\

\strut

\end{minipage} & 2 & \begin{minipage}[t]{\linewidth}\raggedright

จักรินต้องเลือก 1 เมืองเน้นอุตสาหกรรม

ผลรวมของความสุขของทูตจากเมืองอุตสาหกรรมที่เดินทางไปเมืองหลวงจะเท่ากับ 2\\

\strut \\

\includegraphics[width=1.61458in,height=1.64583in,alt={image.png}]{/public/upload/a7c952a7d7.png}\\

\strut

\end{minipage} \\

\begin{minipage}[t]{\linewidth}\raggedright

8 5\\

7 5\\

1 7\\

6 1\\

3 7\\

8 3\\

2 1\\

4 5\\

\strut

\end{minipage} & 9 & \begin{minipage}[t]{\linewidth}\raggedright

\hfill\break

\strut

\end{minipage} \\

\end{longtable}

}



\subsection*{Input}
บรรทัดที่ 1: จำนวนเต็ม n และ k (2 ≤ 𝑛 ≤ 2⋅10\^{}5, 1 ≤ 𝑘 \textless 𝑛) แทน

จำนวนเมืองทั้งหมด และจำนวนเมืองเน้นอุตสาหกรรม ตามลำดับ



บรรทัดที่ 2: เป็นต้นไปจนถึง n-1: จำนวนเต็ม 2 ตัว u และ v (1 ≤ 𝑢, 𝑣 ≤𝑛)

แทนถนนที่เชื่อมต่อเมือง u และ v



\subsection*{Output}
บรรทัดที่ 1: จำนวนเต็ม 1 ตัว แทนผลรวมของความสุขของทูตทั้งหมดที่มากที่สุด\\



\end{document}
